\subsection{What Relational Evidence Cannot Separate}
A group of friends that deliberately routes favours to each other and a ring that
farms ratings produce the same structure, the same money cycles and the same
deliberate routing. Nothing in the friendship graph, the payment graph or the offer
log distinguishes them; only intent does. This is an identifiability limit, not a
tuning problem. The policy test asks whether a group took its own errands more often
than friends elsewhere on campus do. Honest favour-trading answers yes for the same
reason a ring does, and a careful ring that farms only as often as ordinary friends
help each other answers no. Fig.~\ref{fig:sweeps} shows both ends: honest groups are
flagged once favours become frequent, and the deployed test loses coordinated rings
as they farm less. Only the content of the errands or a person can separate the two
cases, and our language channel is not yet reliable enough to lean on
(Section~\ref{sec:eval-system}). We therefore keep a person between every ring flag
and any sanction: ring and farming flags never feed the automatic ladder, and the
automated parts of the system may demote and pause but never exclude.

\subsection{Why the First Policy Test Failed}
The first version of the test used the router's own logged scores as its model of
who would accept an errand. Within a tier every candidate is notified at once, so the
order of the list has little effect on who answers first, and the scores overstated
how much the router decides. A coordinated ring exploited exactly that. When the
partner is always the router's first choice, a test that believes the router expects
the partner to take the errand nearly every time and finds nothing surprising. Three
results in Section~\ref{sec:eval-policy} point to the same cause. Exploration helped
only slowly, setting the trust weight in the expectation to zero restored recall, and
an acceptance model fitted to what runners actually did restored most of it at the
deployed exploration rate, without that cost. The general lesson is that a detector
conditioned on a matcher should condition on the matcher's measured effect, not on
its design weights.

\subsection{Limitations}
Every result here is from simulation. The friendship graph is a block model, not a
real campus, and runner behaviour (willingness, response delay, friend affinity) is
assumed. We varied the parameters that matter most, and the confidence intervals
describe variation across seeds of that model, not uncertainty about whether the
model is right. A pilot may find behaviour we did not imagine. The thresholds in
Table~\ref{tab:params} are reasoned defaults, not values fitted to data; the offer log
exists partly so that a pilot can fit them. The acceptance model has four features,
and a ring that varies distance and timing to look like an ordinary friend will be
harder to separate than the coordinated rings simulated here. A ring whose members
never declare friendship is invisible to every friendship-based mechanism, at the
price of forgoing the trust boost; its cycles would appear only in the any-pair cycle
search, which flags too many honest students to act on. The breadth attack defeats the
reputation weighting at the cost of eighty real accounts. At a shop with few honest
reporters, two or three colluders can move its reference, although the near-line rule
still sees them. A runner can still claim a priced item's money under the name of an
unpriced line on the same order, bounded only by the escrow hold. Line-walkers keep
being paid the allowance until an administrator reviews the flag, so review capacity
is part of the defence. Finally, the language-channel evaluation uses eleven
hand-written cases.

\subsection{Ethics and Deployment}
These mechanisms withhold money and can pause or suspend students, so we built them
around three rules. Automated components may warn, demote and pause, but only a person
may exclude: account suspension needs flags an administrator has upheld. Every flag
carries the evidence that raised it, every sanction tells the student why, and a
dismissed flag releases the money and lifts the pause it caused. And the advisory
model runs on campus hardware, so that students' errand notes and reviews do not
leave the institution. A pilot will collect data about identifiable students and will
need institutional ethics approval and informed consent before it begins.

\subsection{Beyond Campus Errands}
Nothing in the design is specific to food. It applies wherever expenses are reported
rather than documented, such as shared purchasing in hostels or offices and informal
neighbourhood delivery, and more generally to any platform that uses social ties both
to route work and to look for collusion. Once the router shapes who meets whom, the
fraud detector must account for the router, or it ends up auditing the router's
choices instead of people's.
